% ============================================================
%  ch00_karta_tytulowa.tex  —  Front Matter Passages
%  Dedication, epigraphs, and title-page text
% ============================================================

\chapter*{Karta tytułowa}
\addcontentsline{toc}{chapter}{Karta tytułowa}
\markboth{Karta tytułowa}{}

% ── Opening line ─────────────────────────────────────────────
\section{Motto}

\poltext{Zatrute dzieciństwo nie pozwoli zapomnieć.}
\poltrans{A poisoned childhood will not allow one to forget.}

\textbf{Zatrute} — \textit{poisoned} (adjective from \textit{zatruć}, to poison). 
Neuter nominative singular, agreeing with \textit{dzieciństwo}. 
Adjectival agreement exactly as in Latin.

\textbf{dzieciństwo} — \textit{childhood} (neuter noun, nominative singular, subject). 
The suffix \textit{-stwo} is a productive Polish nominaliser, like English 
\textit{-hood} or \textit{-ness}. Root: \textit{dziecko} (child).

\textbf{nie pozwoli} — \textit{will not allow}. Third person singular future 
perfective of \textit{pozwolić}. The perfective aspect signals a definitive, 
completed future action — not an ongoing permission withheld, but an absolute refusal.

\textbf{zapomnieć} — \textit{to forget} (perfective infinitive). 
From \textit{zapominać/zapomnieć}. The \textit{za-} prefix perfectivises 
\textit{pamiętać} (to remember) and reverses its meaning: to have-remembered-away.

\grambox{Impersonal construction}{\textit{nie pozwoli zapomnieć} uses an 
impersonal infinitive — there is no explicit object (one, you, we). This 
gives the line a universal quality. The book's subtitle is 
\textit{Zapomniana epidemia} (the forgotten epidemic); this opening line 
directly inverts that word: the epidemic may be forgotten, but a poisoned 
childhood will not \textit{permit} forgetting.}

% ── Dedication ───────────────────────────────────────────────
\section{Dedication}

\poltext{Ojcu, Hubertowi\\Synowi, Michałowi}
\poltrans{To [my] father, Hubert\\To [my] son, Michał}

Both lines use the \textbf{dative case} — the case of giving and dedicating. 
English requires the preposition ``to''; Polish encodes it in the noun ending.

\polword{Ojciec}{father} → \polword{Ojcu}{to [my] father} (dative singular).

\polword{Syn}{son} → \polword{Synowi}{to [my] son} (dative singular).

Proper names follow exactly the same pattern: 
\polword{Hubert}{} → \polword{Hubertowi}{}; 
\polword{Michał}{} → \polword{Michałowi}{}. 
Even personal names decline fully in Polish — a feature that would feel 
strange in English but is entirely natural here, and is identical to Latin.

The parallel structure — father and son, two generations — mirrors each 
other perfectly in form as well as meaning. No verb, no elaboration. 
The author dedicates a book about the damage done to children to his own 
father and his own son. That juxtaposition is quietly devastating.

% ── Stephen King epigraph ────────────────────────────────────
\section{Epigraph}

\poltext{Czasami to, o czym chcielibyśmy zapomnieć, co miało odejść 
w przeszłość, wcale nie odchodzi.}
\poltrans{Sometimes that which we would like to forget, that which was 
supposed to recede into the past, does not go away at all.}

\textbf{Czasami} — \textit{sometimes}. A gentle, resigned opening.

\textbf{to, o czym\ldots} — \textit{that which\ldots} Relative clause construction. 
\textit{To} is the antecedent; \textit{o czym} is the relative pronoun in the 
locative (literally ``about which'').

\textbf{chcielibyśmy zapomnieć} — \textit{we would like to forget}. 
\textit{Chcielibyśmy} is the conditional first person plural of \textit{chcieć} 
(to want). The \textit{-by-} infix is the Polish conditional marker, equivalent 
to French \textit{-rais/-rait}.

\textbf{co miało odejść} — \textit{that which was supposed to go away}. 
\textit{Miało} is the neuter past of \textit{mieć} used as an auxiliary 
meaning ``was supposed to / was meant to.'' \textit{Odejść}: the prefix 
\textit{od-} (away from) combined with \textit{iść} (to go).

\textbf{w przeszłość} — \textit{into the past}. \textit{Przeszłość} in the 
accusative, required by the motion of \textit{odejść} — going \textit{into} 
somewhere. Root: \textit{przeszły} (past, gone).

\textbf{wcale nie} — \textit{not at all}. \textit{Wcale} is an intensifier 
of negation. Compare \textit{nie odchodzi} (doesn't go away) with 
\textit{wcale nie odchodzi} (doesn't go away at all — simply refuses to leave).

\textbf{odchodzi} — \textit{goes away} (imperfective present). The imperfective 
aspect emphasises the ongoing, continuous failure to leave.

The sentence is structured as a long double relative clause — two layers of 
``that which we wanted gone'' — before the short, flat conclusion: 
\textit{wcale nie odchodzi}. Note the echo of \textit{zapomnieć} from the 
very first line of the book.

\vspace{1em}

\poltext{Czasami to do ciebie wraca.\\{\small --- Stephen King, \textit{To}}}
\poltrans{Sometimes it comes back to you.\\--- Stephen King, \textit{It}}

\textbf{to} — \textit{it} — doing double duty. A neutral demonstrative pronoun 
in Polish, but also the Polish title of King's novel \textit{It}. The sentence 
is simultaneously a quote about unresolvable horror and a quiet nod to the 
book it comes from.

\textbf{do ciebie} — \textit{to you}. \textit{Do} takes the genitive; 
\textit{ciebie} is the emphatic genitive/accusative of \textit{ty} (you). 
The emphatic form rather than the shorter \textit{do cię} gives it 
deliberate, personal weight.

\textbf{wraca} — \textit{returns, comes back}. Imperfective present, third 
person singular. The imperfective aspect: it doesn't just return once, 
it keeps returning.

The author opens with his own sentence about things that refuse to leave, 
then follows with King — whose novel \textit{It} is literally about childhood 
trauma that was supposed to stay buried but keeps returning. The Polish title 
\textit{To} makes the pronoun in the sentence indistinguishable from the 
title itself. After the previous sentence's elegiac tone, this lands with 
quiet menace.

\grambox{Już vs.\ jeszcze — first encounter}{\textbf{Już} marks a transition 
or arrival: \textit{już jest} — he's here now; \textit{już nie} — not anymore. 
\textbf{Jeszcze} marks continuation or addition: \textit{jeszcze jest} — he's 
still here; \textit{jeszcze nie} — not yet. Latin parallel: \textit{już} 
resembles \textit{iam}; \textit{jeszcze} resembles \textit{adhuc}.}
